% OnBuch — Fascicule de Mathématiques Terminale C (compilé avec tectonic/XeLaTeX)
\documentclass[11pt,a4paper,openany]{book}

\usepackage{fontspec}
\usepackage[french]{babel}
\usepackage[a4paper,margin=2.2cm,top=2.6cm,bottom=2.4cm,headheight=14pt]{geometry}
\usepackage{amsmath,amssymb,amsfonts,mathtools}
\usepackage{xcolor}
\usepackage[most]{tcolorbox}
\usepackage{enumitem}
\usepackage{fancyhdr}
\usepackage{titlesec}
\usepackage{array,booktabs}
\usepackage{multicol}
\usepackage{pgffor}
\usepackage{etoolbox}
\usepackage{tikz}
\usetikzlibrary{arrows.meta,positioning,calc,decorations.pathreplacing,angles,quotes}
\usepackage{pgfplots}
\pgfplotsset{compat=1.18}
\usepackage[hidelinks]{hyperref}
\hypersetup{bookmarksnumbered}
\usepackage{pdfpages}
\setcounter{tocdepth}{1}   % sommaire : chapitres + sections uniquement

% ── Palette OnBuch ──────────────────────────────────────────────
\definecolor{onbo}{HTML}{F2620E}
\definecolor{onbo2}{HTML}{FF8A2B}
\definecolor{onbgreen}{HTML}{1E9E63}
\definecolor{onbink}{HTML}{241B12}
\definecolor{onbpaper}{HTML}{FFFBF7}
\definecolor{onbsoft}{HTML}{FFF1E3}
\definecolor{onbline}{HTML}{E8DDD3}
\definecolor{onbmuted}{HTML}{8A7B6D}
\definecolor{onbblue}{HTML}{2D6CDF}

\setlist[enumerate]{leftmargin=1.7em,itemsep=3pt,topsep=3pt}
\setlist[itemize]{leftmargin=1.5em,itemsep=2pt,topsep=2pt}

% ── En-têtes / pieds ────────────────────────────────────────────
\pagestyle{fancy}
\fancyhf{}
\renewcommand{\headrulewidth}{0pt}
\fancyhead[LE]{\small\textsf{\color{onbmuted}\thepage\quad·\quad OnBuch}}
\fancyhead[RE]{}
\fancyhead[LO]{}
\fancyhead[RO]{\small\textsf{\color{onbmuted}\leftmark\quad·\quad\thepage}}
\fancyfoot[C]{\small\textsf{\color{onbmuted}Mathématiques · Terminale C · Baccalauréat}}
\renewcommand{\chaptermark}[1]{\markboth{#1}{}}

% ── Titres de chapitre (gros numéro orange) ─────────────────────
\titleformat{\chapter}[display]
  {\sffamily}
  {\hfill\fontsize{72}{72}\selectfont\bfseries\color{onbo}\thechapter}
  {6pt}
  {\Huge\bfseries\color{onbink}\filright}
  [\vspace{4pt}{\color{onbo}\titlerule[2.5pt]}]
\titlespacing*{\chapter}{0pt}{10pt}{26pt}

\titleformat{\section}
  {\sffamily\Large\bfseries\color{onbink}}{\color{onbo}\thesection}{0.6em}{}
  [{\color{onbline}\titlerule[1pt]}]
\titleformat{\subsection}
  {\sffamily\large\bfseries\color{onbo}}{\thesubsection}{0.6em}{}

% ── Compteurs d'exercices par chapitre ──────────────────────────
\newcounter{exo}[chapter]
\newcommand{\exo}[1]{\par\medskip\noindent{\color{onbo}\sffamily\bfseries\large Exercice \refstepcounter{exo}\theexo}\;{\sffamily\bfseries #1}\par\nopagebreak\smallskip}
\newcommand{\methode}[1]{\par\medskip\noindent{\color{onbblue}\sffamily\bfseries Méthode\;:\;}{\sffamily\bfseries\color{onbink}#1}\par\smallskip}

% Rubrique d'exercices (catégorie thématique)
\newcommand{\rubrique}[1]{\par\medskip\noindent
  {\color{onbo}\sffamily\bfseries\large$\blacktriangleright$\;#1}%
  \par\vspace{-4pt}\noindent{\color{onbline}\rule{\linewidth}{0.8pt}}\par\smallskip}

% Niveau de difficulté (étoiles) : \diff{1..3}
\newcommand{\diff}[1]{\hfill{\color{onbo2}\small%
  \ifcase#1\or\faStarSolid\or\faStarSolid\faStarSolid\or\faStarSolid\faStarSolid\faStarSolid\fi}}
% étoile simple (sans fontawesome) pour robustesse
\renewcommand{\diff}[1]{{\color{onbo2}\small\ifcase#1\or$\star$\or$\star\star$\or$\star\star\star$\fi}}

% Légende de figure centrée
\newcommand{\fig}[1]{\par\vspace{2pt}\centerline{\footnotesize\itshape\color{onbmuted}#1}\par\smallskip}
% Encadré "figure/illustration"
\newtcolorbox{illus}{colback=white,colframe=onbline,boxrule=0.8pt,arc=4pt,
  left=8pt,right=8pt,top=8pt,bottom=6pt,halign=center,breakable}

% ── Boîtes pédagogiques ─────────────────────────────────────────
\newtcolorbox{definition}[1][]{colback=onbpaper,colframe=onbo,boxrule=0pt,
  leftrule=3pt,arc=3pt,left=12pt,right=12pt,top=8pt,bottom=8pt,breakable,
  title={\sffamily\bfseries Définition},coltitle=onbo,colbacktitle=onbpaper,
  fonttitle=\sffamily\bfseries,attach title to upper={\quad#1\par\smallskip}}

\newtcolorbox{theoreme}[1][]{colback=white,colframe=onbgreen,boxrule=0pt,
  leftrule=3pt,arc=3pt,left=12pt,right=12pt,top=8pt,bottom=8pt,breakable,
  title={\sffamily\bfseries Théorème},coltitle=onbgreen,colbacktitle=white,
  fonttitle=\sffamily\bfseries,attach title to upper={\;\;\textmd{\itshape#1}\par\smallskip}}

\newtcolorbox{propriete}[1][]{colback=white,colframe=onbgreen!75,boxrule=0pt,
  leftrule=3pt,arc=3pt,left=12pt,right=12pt,top=8pt,bottom=8pt,breakable,
  title={\sffamily\bfseries Propriété},coltitle=onbgreen!85!black,colbacktitle=white,
  fonttitle=\sffamily\bfseries,attach title to upper={\;\;\textmd{\itshape#1}\par\smallskip}}

\newtcolorbox{remarque}{colback=onbsoft,colframe=onbsoft,boxrule=0pt,arc=3pt,
  left=12pt,right=12pt,top=6pt,bottom=6pt,breakable,
  before upper={\textbf{\color{onbo}Remarque.}\ }}

\newtcolorbox{exemple}{colback=white,colframe=onbline,boxrule=0.8pt,arc=3pt,
  left=12pt,right=12pt,top=6pt,bottom=6pt,breakable,
  before upper={\textbf{\color{onbink}\sffamily Exemple.}\ }}

\newtcolorbox{aretenir}{enhanced,colback=onbsoft,colframe=onbo,boxrule=1pt,arc=6pt,
  left=12pt,right=12pt,top=8pt,bottom=8pt,breakable,
  title={\sffamily\bfseries À retenir},coltitle=white,colbacktitle=onbo,fonttitle=\sffamily\bfseries}

% ── Synthèse de fin de chapitre (« L'essentiel — fiche de révision ») ──
\newtcolorbox{synthese}[1][L'essentiel du chapitre]{enhanced,breakable,
  colback=onbink,colframe=onbink,boxrule=0pt,arc=6pt,
  left=14pt,right=14pt,top=10pt,bottom=10pt,
  fontupper=\color{white},
  title={\sffamily\bfseries\large #1},coltitle=white,colbacktitle=onbo,
  fonttitle=\sffamily\bfseries}

% Sujet « type Baccalauréat » (mise en avant)
\newcounter{sujet}[chapter]
\newcommand{\sujetbac}[1][]{\par\medskip\noindent
  {\color{onbo}\sffamily\bfseries\large Sujet \refstepcounter{sujet}\thesujet\ — type Baccalauréat}%
  \ifstrempty{#1}{}{\;{\sffamily\itshape\color{onbmuted}(#1)}}%
  \par\vspace{-4pt}\noindent{\color{onbo}\rule{\linewidth}{1.2pt}}\par\smallskip}

% Correction (titre vert, sécable)
\newcommand{\corrige}[1]{\par\smallskip\noindent{\color{onbgreen}\sffamily\bfseries\small$\blacktriangleright$ Corrigé #1}\par\nopagebreak\smallskip}

% Raccourcis maths
\newcommand{\R}{\mathbb{R}}
\newcommand{\N}{\mathbb{N}}
\newcommand{\Z}{\mathbb{Z}}
\newcommand{\Q}{\mathbb{Q}}
\newcommand{\C}{\mathbb{C}}
\newcommand{\dd}{\,\mathrm{d}}
\DeclareMathOperator{\Card}{Card}
